\documentclass[12pt,a4paper]{article}
\usepackage[utf8]{inputenc}
\usepackage[T1]{fontenc}
\usepackage{amsmath,amssymb,amsthm,mathrsfs,mathtools}
\usepackage[margin=1in]{geometry}
\usepackage{hyperref}
\hypersetup{colorlinks=true,linkcolor=blue,citecolor=blue,urlcolor=blue}
\usepackage{booktabs,enumitem}

\theoremstyle{plain}
\newtheorem{theorem}{Theorem}[section]
\newtheorem{proposition}[theorem]{Proposition}
\newtheorem{lemma}[theorem]{Lemma}
\theoremstyle{definition}
\newtheorem{definition}[theorem]{Definition}
\theoremstyle{remark}
\newtheorem{remark}[theorem]{Remark}

\newcommand{\R}{\mathbb{R}}
\newcommand{\Z}{\mathbb{Z}}
\newcommand{\C}{\mathbb{C}}
\newcommand{\dd}{\mathrm{d}}
\newcommand{\mm}{\mathfrak{m}}
\newcommand{\cM}{\mathcal{M}}
\DeclareMathOperator{\Res}{Res}

\title{Modular Thermality of the Celestial Spectral Weight:\\[4pt]
Planck Form, Matsubara Poles, and the Moment Family}

\author{Daniel Toupin\\[4pt]
\textit{Golden Physics Project}\\[2pt]
\texttt{goldenphysics.org}}

\date{August 2026 --- v2, superseding \texttt{blackbody\_law\_qg\_dtoupin\_v1}}

\begin{document}
\maketitle

\begin{abstract}
\noindent
The weight $P(\lambda)=\pi\lambda/\sinh(\pi\lambda)$ governs boost
integrals throughout celestial holography. A companion paper derives it
from two-particle phase space, where it appears as
$\Gamma(1+i\lambda)\Gamma(1-i\lambda)$ on the shadow locus. This paper
collects its thermal and modular structure, with every identity
recomputed independently to at least 24 digits.

The weight takes a Planck form,
$P=2\pi\lambda[n_B(\pi\lambda)-n_B(2\pi\lambda)]$, in which the
zero-point energies of the two Bose terms cancel exactly. Read through
the Bisognano--Wichmann theorem, it is the partition function of one
oscillator at the modular temperature $\beta=2\pi$, weighted by boost
frequency. Five further descriptions of the same function are given:
squared Gamma modulus on the critical line, characteristic function of
the logistic distribution, Weierstrass product over the Laplace
spectrum of the thermal circle $\R/2\pi\Z$, meromorphic function whose
poles sit exactly at the bosonic Matsubara points $\lambda=in$ with
residues $i(-1)^n n$, and generator of cumulants that are even zeta
values. The Mellin moments form a one-parameter family,
\[
\mm_s=\frac{1}{2\pi}\int_0^\infty\lambda^{s-1}P(\lambda)\dd\lambda
=\pi^{-(s+1)}\bigl(1-2^{-(s+1)}\bigr)\Gamma(s+1)\zeta(s+1),
\]
verified at five values of $s$; the members $\mm_1=1/8$ and
$\mm_3=1/16$ are exact, and the family is the structural twin of the
Stefan--Boltzmann integral, with $\zeta(s+1)$ where the cavity law has
$\zeta(4)$.

The physical thesis of the superseded draft is withdrawn. The claim
that cavity radiation and the ultraviolet behaviour of quantum gravity
are one disease with one cure rested on a companion claim, since
retracted, that $P$ replaces the Feynman loop measure. It does not.
What remains is a precise thermal reading of a function that arises
from phase space, which is worth stating on its own terms and is not
worth more than that.
\end{abstract}

\medskip
\noindent\textbf{Keywords:} celestial holography, modular theory,
Bisognano--Wichmann, Matsubara frequencies, Dirichlet lambda function.

\tableofcontents
\newpage

% =====================================================================
\section{What the weight is}
% =====================================================================

The superseded draft opened by positing $P(\lambda)$ as the measure
that cures loop divergences, then sought a thermal reading of it. The
present order is the reverse, and it is the order the evidence
supports.

\begin{theorem}[Origin, from \cite{ToupinLoops}]
\label{thm:origin}
The Mellin image of two-particle massless phase space in celestial
coordinates is
\begin{equation}
\int\dd\Pi_2\,\omega_5^{\Delta_5-1}\omega_6^{\Delta_6-1}
=\frac{1}{8\pi}\Bigl(\frac{M}{2}\Bigr)^{\Delta_5+\Delta_6-2}
\frac{\Gamma(\Delta_5)\Gamma(\Delta_6)}{\Gamma(\Delta_5+\Delta_6)} ,
\end{equation}
and on the shadow locus $\Delta_5+\Delta_6=2$ with
$\Delta_5=1+i\lambda$ the Gamma ratio equals
\begin{equation}
\Gamma(1+i\lambda)\Gamma(1-i\lambda)=\frac{\pi\lambda}{\sinh(\pi\lambda)}
=:P(\lambda).
\label{eq:P}
\end{equation}
\end{theorem}

So $P$ is a phase-space density of a unitarity cut. Everything below
describes that density; nothing below licenses replacing a measure
anywhere.

% =====================================================================
\section{Planck form and the modular temperature}
% =====================================================================

\begin{theorem}[Planck form, with cancelling zero-point energy]
\label{thm:planck}
Write $n_B(y)=(e^y-1)^{-1}$ and $E(x)=\tfrac{x}{2}\coth\tfrac{x}{2}$.
For all $\lambda>0$,
\begin{equation}
P(\lambda)
=2\pi\lambda\bigl[n_B(\pi\lambda)-n_B(2\pi\lambda)\bigr]
=2\pi\lambda\sum_{n\ge0}e^{-2\pi\lambda(n+\frac12)},
\qquad
2E(\pi\lambda)-E(2\pi\lambda)=P(\lambda).
\label{eq:planck}
\end{equation}
The zero-point contributions cancel identically:
$2\cdot\tfrac{\pi\lambda}{2}-\tfrac{2\pi\lambda}{2}=0$.
\end{theorem}

\begin{proof}
Both reduce to $\coth y-\coth2y=1/\sinh2y$. Verified at
$\lambda=0.4,1.3,3.1$ (first form) and $\lambda=0.9,2.6$ (second) to
relative error below $3.3\times10^{-39}$.
\end{proof}

\begin{theorem}[Gibbs reading at $\beta=2\pi$]
\label{thm:gibbs}
Let $K$ generate the Lorentz boost preserving a Rindler wedge whose
edge is the celestial sphere. By the Bisognano--Wichmann theorem
\cite{BW1975,BW1976}, the Minkowski vacuum restricted to the wedge
algebra is a KMS state at inverse temperature $\beta=2\pi$ in boost
time. The middle expression of \eqref{eq:planck} is the partition
function of a single oscillator with level spacing $2\pi\lambda$ in
that state, multiplied by the boost frequency $2\pi\lambda$.
\end{theorem}

The cancellation in \eqref{eq:planck} deserves a word. A single Bose
mode carries $\tfrac{x}{2}$ of zero-point energy, which is what makes
the naive sum over modes diverge. The combination appearing in $P$ is a
difference of two such modes at spacings $\pi\lambda$ and $2\pi\lambda$,
and the zero-point pieces cancel between them. The finiteness of $P$ at
large $\lambda$ therefore has a clean thermodynamic reading. It is a
statement about this function, and it says nothing about loop
integration; see Sec.~\ref{sec:retract}.

% =====================================================================
\section{Five further faces of the same function}
% =====================================================================

\begin{proposition}[Equivalent descriptions]
\label{prop:faces}
Let $P(\lambda)=\pi\lambda/\sinh(\pi\lambda)$.
\begin{enumerate}[label=\textup{(\roman*)}]
\item \emph{Gamma modulus on the critical line.}
$P(\lambda)=\Gamma(1+i\lambda)\Gamma(1-i\lambda)=|\Gamma(2s)|^2$ at
$s=\tfrac12+\tfrac{i\lambda}{2}$. Checked to $7\times10^{-41}$.
\item \emph{Logistic characteristic function.}
$P(\lambda)=\int_\R e^{i\lambda x}\,\dd x/\bigl(4\cosh^2(x/2)\bigr)$,
so $P$ is the characteristic function of the logistic distribution,
equivalently the Fourier transform of the reflectionless
P\"oschl--Teller ground-state density; the transform in the other
direction is $\widehat P(k)=\tfrac{\pi}{2}\operatorname{sech}^2(k/2)$.
Checked to $1.5\times10^{-31}$.
\item \emph{Functional determinant.}
$P(\lambda)^{-1}=\prod_{n\ge1}\bigl(1+\lambda^2/n^2\bigr)$, a
Weierstrass product over the Laplace spectrum of the thermal circle
$\R/2\pi\Z$. Checked at $\lambda=0.8,2.5$ to $3.1\times10^{-31}$.
\item \emph{Matsubara poles.}
The continuation of $P$ is meromorphic with simple poles exactly at
$\lambda=in$, $n\in\Z\setminus\{0\}$, and
$\Res_{\lambda=in}P(\lambda)=i(-1)^n n$. Checked at $n=1,2,3$ to
$10^{-12}$ by contour evaluation. At $\beta=2\pi$ the bosonic Matsubara
frequencies are $\omega_n=2\pi n/\beta=n$, so the poles sit at the
Matsubara points of the modular temperature.
\item \emph{Cumulants are even zeta values.}
For $|\lambda|<1$,
$\log P(\lambda)=-\sum_{k\ge1}\frac{(-1)^{k+1}\zeta(2k)}{k}\lambda^{2k}$.
\end{enumerate}
\end{proposition}

Item (iv) is the sharpest of the five for physical purposes. A weight
whose poles land on Matsubara frequencies at $\beta=2\pi$ behaves,
under contour deformation, the way a thermal propagator behaves at that
temperature. That is a structural fact about boost integrals in flat
space, and it is the reason the same function keeps appearing in
otherwise unrelated celestial calculations.

% =====================================================================
\section{The moment family}
% =====================================================================

\begin{theorem}[Moment family]
\label{thm:sb}
For every real $s>0$,
\begin{equation}
\mm_s:=\frac{1}{2\pi}\int_0^\infty\lambda^{s-1}P(\lambda)\dd\lambda
=\pi^{-(s+1)}\bigl(1-2^{-(s+1)}\bigr)\Gamma(s+1)\,\zeta(s+1).
\label{eq:sb}
\end{equation}
In particular $\mm_1=\tfrac18$ and $\mm_3=\tfrac{1}{16}$.
\end{theorem}

\begin{proof}
Insert the exponential series of \eqref{eq:planck}. Termwise
integration is justified by monotone convergence, since every term is
positive:
\begin{equation}
\mm_s=\int_0^\infty\lambda^s\sum_{k\ge0}e^{-(2k+1)\pi\lambda}\dd\lambda
=\frac{\Gamma(s+1)}{\pi^{s+1}}\sum_{k\ge0}(2k+1)^{-(s+1)},
\end{equation}
and the odd-integer sum is the Dirichlet lambda function
$(1-2^{-w})\zeta(w)$ at $w=s+1$. Verified independently at
$s=0.7,1,2,3,4.5$ with relative error at most $2.5\times10^{-31}$; the
values $\mm_1=0.125$ and $\mm_3=0.0625$ came out exact at working
precision.
\end{proof}

\begin{remark}[The analogy, sized correctly]
The Stefan--Boltzmann law arises from $\int_0^\infty
x^3/(e^x-1)\dd x=\Gamma(4)\zeta(4)=\pi^4/15$. Equation \eqref{eq:sb}
has the same shape, with the Dirichlet lambda function in place of
zeta, because the Bose difference in \eqref{eq:planck} restricts the
sum to odd integers. The two formulas are cousins in the sense that
both are Mellin transforms of Bose-type kernels. The superseded draft
went further and treated the resemblance as evidence that one law cures
the pathology of the other. It is not evidence of that, and the
inference is withdrawn.
\end{remark}

\begin{remark}[On $\mm_1=\cM_1=1/8$]
The ladder moment $\cM_1$ of \cite{ToupinSpectral} is the $s=1$ member
of this family, so those two $1/8$s agree by construction. A third
$1/8$, appearing in the digamma moment of \cite{ToupinBlocks}, comes
from an integral over the full line and is a separate evaluation. The
agreement of that one with these remains unexplained, and should be
read as a coincidence until someone derives it.
\end{remark}

% =====================================================================
\section{Retractions}
\label{sec:retract}
% =====================================================================

The following are withdrawn from
\texttt{blackbody\_law\_qg\_dtoupin\_v1}, and are listed by name so
that readers holding that draft can mark it.

\begin{enumerate}
\item \textbf{``The ultraviolet catastrophe of cavity radiation and the
ultraviolet divergence of quantum gravity are the same disease, and
admit the same cure.''} Withdrawn. The claim depended on $P$ acting as
a replacement for the Feynman loop measure, which is retracted in
\cite{ToupinSpectral}. The companion derivation \cite{ToupinLoops}
reconstructs the standard loop integral from its cut exactly,
divergences included.

\item \textbf{``$P$ renders the celestial loop expansion finite at
every order.''} Withdrawn. The bound $\cM_L\le(1/8)^L$ is true of the
convolution integrals in \cite{ToupinSpectral} and carries no
implication for amplitudes.

\item \textbf{``$\mm_1=1/8$ is the gravitational Stefan--Boltzmann
constant.''} The value is correct, and Theorem \ref{thm:sb} proves it.
The title attached to it presumes the loop-measure reading, so the name
is withdrawn while the number stands.

\item \textbf{Claims inherited from the companion drafts.} The
description of $P$ as the Plancherel density of $\mathrm{SL}(2,\C)$ is
false; the density for the scalar principal series grows like
$\lambda^2$. The fifteen-digit box verification cited in support was
circular, since the routine it used reparameterizes a Feynman integral
and compares it with itself.
\end{enumerate}

% =====================================================================
\section{Status of each statement}
% =====================================================================

\begin{center}
\begin{tabular}{lll}
\toprule
Statement & Status & Check \\
\midrule
$P$ from phase space & proved in \cite{ToupinLoops} & $10^{-41}$ \\
Planck form; zero-point cancellation & proved & $3.3\times10^{-39}$ \\
Gibbs reading at $\beta=2\pi$ & exact identification & \\
Gamma modulus on critical line & proved & $7\times10^{-41}$ \\
Logistic characteristic function & proved & $1.5\times10^{-31}$ \\
Weierstrass product & proved & $3.1\times10^{-31}$ \\
Matsubara poles, residues $i(-1)^nn$ & proved & $10^{-12}$ \\
Cumulants are even zeta values & classical & \\
Moment family \eqref{eq:sb} & proved & $2.5\times10^{-31}$ \\
$\mm_1=1/8$, $\mm_3=1/16$ & proved & exact \\
Cavity radiation as a cure for QG & \textbf{withdrawn} & \\
Finiteness of the loop expansion & \textbf{withdrawn} & \\
\bottomrule
\end{tabular}
\end{center}

\noindent
All numerical checks are reproduced by \texttt{verify\_blocks\_v2.py},
archived with this paper.

\begin{thebibliography}{9}
\bibitem{ToupinLoops} D.~Toupin, ``Loop amplitudes from unitarity cuts
in celestial Mellin space: a complete derivation for the scalar box,''
v19, 2026.
\bibitem{ToupinSpectral} D.~Toupin, ``The spectral weight
$\pi\lambda/\sinh\pi\lambda$: derivation from phase space and exact
moments,'' v2, 2026.
\bibitem{ToupinBlocks} D.~Toupin, ``Principal-series blocks on the
celestial sphere,'' v2, 2026.
\bibitem{BW1975} J.~Bisognano, E.~Wichmann, J.\ Math.\ Phys.\ 16 (1975)
985.
\bibitem{BW1976} J.~Bisognano, E.~Wichmann, J.\ Math.\ Phys.\ 17 (1976)
303.
\end{thebibliography}

\end{document}
